% PDF-derived LaTeX transcription. Original manuscript source was unavailable.
% Published PDF SHA-256: 98b5b27331cf697b5f4253aea7e37ce27848d98dc658aede7be38f1dc55983fd
% Extracted text SHA-256: d3da47c96a33247bf5f828165efef347279e216f1564bd6ad6225a580d655633
\documentclass[11pt]{article}
\usepackage[margin=1in]{geometry}
\usepackage{amsmath,amssymb}
\setlength{\emergencystretch}{3em}
\title{Energy Balance and Comparison Observables in a PMMA Impact-Fracture Simulation Audit}
\author{}
\date{}
\begin{document}
\maketitle
\noindent\textbf{PDF-derived LaTeX transcription.} This editable text follows the PDF extraction. Refer to the published PDF for figures, table layouts, and mathematical notation.

% Extracted chunk 1

Energy Balance and Comparison Observables in a PMMA
Impact-Fracture Simulation Audit

\section*{Abstract}

We examine retained computations for two PMMA impact specimens using an affine particle-
in-cell/material point method coupled to irreversible AT2 damage. The question is whether
the recorded energy and field comparisons support treating the exploratory coupled setup as
a predictive fracture model. For TAF1, the declared relative energy residual decreases from
88.3429\% to 6.9786\% over five retained values of the mechanics Courant–Friedrichs–Lewy (CFL)
parameter. This is a substantial reduction, but leaves a finite defect and crosses a source
change that exposed the CFL control. On seven common comparison times through 29.976 \ensuremath{\mu}s,
all-node relative displacement discrepancies are 41.9292\% and 64.1526\% for TAF1 and TAF2.
Subtracting the first measured displacement changes these to 42.0446\% and 54.0514\%, with
unchanged velocity discrepancies of 36.4432\% and 52.6399\%. The reference change alters the
comparison, rather than demonstrating a better simulator. First particle damage crossings
of 10\ensuremath{-}3 are distinguished from experimentally resolved crack opening; the homogeneous AT2
relation explains why a small diffuse-damage crossing need not indicate a crack. Stored-array and
component-ledger reductions support these descriptive results. Initial-state, clock, projection
and discretization limits prevent attributing the mismatch uniquely to a constitutive law. The
evidence diagnoses the adequacy of this particular calibration setup and does not establish a
validated high-rate PMMA predictor.

\section*{1 Question and experimental context}

An impact-fracture computation can disagree with measurements through its material response,
boundary history, initial state, discretization or definition of the compared observable. These
possibilities matter in the retained PMMA computations studied here: an apparent early fracture
event, a large energy defect and large field discrepancies were reported together. We ask which
conclusions survive when the energy ledger, displacement reference and damage event are stated
explicitly. The contribution is a quantitative audit of one coupled setup, denoted M0, using its
saved results. It is not a new fracture law or a new particle method.
The experimental reference is the impact study of Eid, Seghir and Réthoré [ 1], with accompa-
nying data [ 2]. Their specimens have an in-plane size of 140 \ensuremath{\times} 70 mm, a circular hole of diameter
30 mm and a 15 mm notch extending from the hole edge. A projectile loads a waveguide attached
to the specimen. The hole converts the compressive input into tensile loading, and reflections at
the free boundary contribute to the subsequent response. The reported projectile speeds for TAF1
and TAF2 are 30.5 and 22 m s\ensuremath{-}1. These are projectile speeds before waveguide impact; they are
not uniform initial specimen velocities.
The reference uses finite-element digital image correlation (DIC), with triangular linear elements
and a crack-conforming twin-node description. Displacements are spatially regularized and filtered
in time using a five-frame, second-order Savitzky–Golay filter before temporal differentiation to
velocity [1]. Thus adjacent field samples are not independent measurements. The published effective
elastic constants were identified from the pre-crack response under a plane-strain assumption using
measured edge displacement histories. Their use below does not independently identify a dynamic
fracture constitutive law.
1

The present analysis uses the two calibration specimens TAF1 and TAF2. Both the prescribed
impacted-edge kinematics and the interior comparison targets derive from the same DIC records.
Interior discrepancies therefore diagnose the calibration setup; they do not constitute a held-out
validation. All computational results below come from retained arrays, reports and source descrip-
tions. No new mechanical simulation or experiment is introduced.

\section*{2 Implemented model and retained comparison}


\section*{2.1 Exploratory material hypothesis}

M0 combines quadratic-kernel affine particle-in-cell (APIC) transfer with a material point method
(MPM) mechanics update and history-driven irreversible AT2 damage. Let sa be the singular values
of the deformation gradient, ea = log(max( sa, 10\ensuremath{-}12)), and \ensuremath{\langle}x\ensuremath{\rangle}+ = max( x, 0), \ensuremath{\langle}x\ensuremath{\rangle}\ensuremath{-} = min( x, 0).
The implemented spectral quadratic-Hencky split and degradation are
\ensuremath{\psi}\ensuremath{\pm} = \ensuremath{\mu}
3\ensuremath{\sum}
a=1
\ensuremath{\langle}ea\ensuremath{\rangle}2
\ensuremath{\pm} + \ensuremath{\lambda}
2
\ensuremath{\langle} 3\ensuremath{\sum}
a=1
ea
\ensuremath{\rangle}2
\ensuremath{\pm}
, g(d) = (1 \ensuremath{-} k)(1 \ensuremath{-} d)2 + k, (1)
\ensuremath{\psi}(F, d) = g(d)\ensuremath{\psi}+(F ) + \ensuremath{\psi}\ensuremath{-}(F ), H(t) = max
\ensuremath{\tau} \ensuremath{\leq}t
\ensuremath{\psi}+(F (\ensuremath{\tau} )). (2)
Here d \ensuremath{\in} [0, 1] is damage, \ensuremath{\lambda} and \ensuremath{\mu} are the Lamé parameters associated with E and \ensuremath{\nu}, and H
is the tensile history. The principal stress is formed from the degraded tensile and undegraded
compressive parts. This describes the implemented hypothesis, without asserting that this finite-
strain extension has been experimentally identified for PMMA.
Table 1: Exploratory M0 parameters and fixed numerical setting. The extrusion is a numerical
dimension, not an established experimental thickness.
Quantity Value
Density \ensuremath{\rho} 1200 kg m\ensuremath{-}3
Young’s modulus E; Poisson ratio \ensuremath{\nu} 6 GPa; 0.3
Fracture energy Gc 417.9326667 J m\ensuremath{-}2
Damage length \ensuremath{\ell}; residual stiffness k 2 mm; 10\ensuremath{-}6
Mechanics grid spacing h 1 mm
Grid dimensions 145 \ensuremath{\times} 75 \ensuremath{\times} 6 nodes
Numerical extrusion bn 2 mm
Particle counts: TAF1; TAF2 142956; 144108
Table 1 specifies the material card used by the retained runs. The value of Gc in Supplement S1
was formed as (1\ensuremath{-}\ensuremath{\nu}2)K2
IC /E using the fitted crack-propagation-law parameter KIC = 1.66 MPa\ensuremath{\surd}m
reported by Eid et al. [ 1] as an exploratory card input. The 2 mm length is an exploratory choice
rather than an independently identified length scale. M0 contains no rate dependence. Neither
the derivation of Gc nor the effective elastic fit establishes that these parameters form a validated
dynamic PMMA law.
On the active reference-grid graph, the recorded damage functional has the form
Ah(d; H) =
\ensuremath{\sum}
i
Vi
[
g(di)Hi + Gc
2\ensuremath{\ell} d2
i
]
+ Gc\ensuremath{\ell}
2h2
\ensuremath{\sum}
\{i,j\}
Vi + Vj
2 (dj \ensuremath{-} di)2. (3)
2

The edge sum counts each active Cartesian nearest-neighbour edge once; Vi is the projected reference
volume. The irreversible lower obstacle is the maximum of the particle-projected lower bound
and the previously converged reference-grid damage. A projected red-black Gauss–Seidel solve
with relaxation factor one updates this field, followed by mapping to particles with non-healing
controls. Mechanics and particle damage storage use float32; the reference-grid damage operator
and projected Karush–Kuhn–Tucker residual use float64. The stopping control is a maximum
projected residual of 10 Pa, with at most 5000 full sweeps. It is a nonlinear solve control, not an
experimental accuracy tolerance. Both baseline reports record convergence of their damage solves.

\section*{2.2 Boundary, time and spatial support}

Measured impacted-edge histories prescribe grid velocity. The source samples this boundary history
at mechanics substep midpoints and separately projects the out-of-plane grid velocity to zero. We
call this the implemented plane-strain constraint, while recognizing that a finite extrusion with that
projection still requires checks of its fields and reactions. The numerical extrusion does not supply
the experimental specimen thickness. Absolute three-dimensional experimental force or energy
cannot be inferred from it.
Seven common physical targets are retained: 0, 4.996, 9.992, 14.988, 19.984, 24.980 and
29.976 \ensuremath{\mu}s. TAF1 uses exact measured samples at native indices 0, 2, 4, 6, 8, 10, 12. TAF2 uses
its exact first sample and piecewise-linear interpolation between adjacent native samples for the
later six targets. There is no additional retiming in this comparison. Integer mechanics substeps
land at each declared target; the CFL parameter controls the schedule while the observation times
remain fixed.
The saved field comparison averages the four equal-in-plane-coordinate particle layers, then
projects to DIC reference nodes using local linear moving weighted least squares. Its compact-
support weight is (1 \ensuremath{-} r/R)2 for r < R , with R = 3 a and particle spacing a. The prescribed
checks require at least six neighbours, design rank three and a normal-matrix condition number
no greater than 108; the lattice-rounding tolerance is 2 \ensuremath{\times} 10\ensuremath{-}5a. Failure of a support or condition
check prohibits extrapolation. These projection settings are shared by the absolute and incremental
comparison contracts. The retained prediction and target arrays have exactly matching coordinates
and times. This analysis reproduces discrepancies of those saved projected arrays, rather than
independently replaying the particle-to-DIC projection.
The all-supported masks contain 3429 nodes for TAF1 and 3465 for TAF2. A separate interior
mask requires x \ensuremath{-} xmin \ensuremath{\geq} 2h, where xmin is the impacted-edge coordinate of the DIC nodes. It
contains 3360 and 3398 nodes. Excluding this boundary band reduces direct inclusion of prescribed-
edge values; it does not make the remaining measurements independent of the imposed boundary
record.

\section*{2.3 Field definitions and normalization}

Let qs
ni and qm
ni denote the saved simulated and measured vectors at time n and node i in a specified
mask M. We report
\ensuremath{\varepsilon}q =
[\ensuremath{\sum}6
n=0
\ensuremath{\sum}
i\ensuremath{\in}M \ensuremath{\parallel}qs
ni \ensuremath{-} qm
ni\ensuremath{\parallel}2
]1/2
[\ensuremath{\sum}6
n=0
\ensuremath{\sum}
i\ensuremath{\in}M \ensuremath{\parallel}qm
ni\ensuremath{\parallel}2
]1/2 , RMSEq =
[ \ensuremath{\sum}6
n=0
\ensuremath{\sum}
i\ensuremath{\in}M \ensuremath{\parallel}qs
ni \ensuremath{-} qm
ni\ensuremath{\parallel}2
7|M|
] 1/2
. (4)
These are uniform node/frame Frobenius and vector root-mean-square discrepancies, without area
or temporal quadrature weights. The three stored vector components are included; the compared
3

out-of-plane components are zero. A global relative discrepancy uses the global target norm, rather
than averaging seven framewise percentages.
Two measured displacement targets are kept distinct:
um,abs
ni = um
ni, u m,inc
ni = um
ni \ensuremath{-} um
0i. (5)
The saved simulation starts with zero displacement; the same simulated displacement is used against
each target. The two target versions have identical underlying displacement samples, geometry,
times and velocity. They differ exactly by subtraction of the first measured displacement. For
the incremental first frame the measured norm and error are both zero, so its framewise relative
discrepancy is undefined; this does not prevent the global norm in Eq. ( 4) from being defined.

\section*{3 Declared energy balance}

The energy diagnostic uses tracked energy E = K + U + \ensuremath{\Gamma}, comprising particle translational plus
APIC affine kinetic energy, recoverable elastic energy and the recorded fracture functional contri-
bution. Let Wg, Wc, Wb be cumulative gravity, contact and prescribed-boundary work. Let \ensuremath{\Delta}Kc,
\ensuremath{\Delta}Kb and \ensuremath{\Delta}Kp be cumulative kinetic changes from contact, boundary and plane-strain projections.
The source declares
Dc = Wc \ensuremath{-} \ensuremath{\Delta}Kc, D b = Wb \ensuremath{-} \ensuremath{\Delta}Kb, D p = \ensuremath{-}\ensuremath{\Delta}Kp, (6)
RE = E + Dc + Db + Dp \ensuremath{-} E0 \ensuremath{-} Wg \ensuremath{-} Wc \ensuremath{-} Wb, (7)
\ensuremath{\eta}E = |RE|
max\{|E0|, |E|, |Wg|, |Wc|, |Wb|, |Dc|, |Db|, |Dp|, 10\ensuremath{-}30 J\} . (8)
The scale in Eq. ( 8) is the implemented numerical ledger scale. It is not simply boundary work,
initial energy or a measured experimental energy. The symbols Dc, Db, Dp preserve the source’s
bookkeeping convention; in particular boundary and plane-strain projection terms are numerical
constraint terms, not measured material dissipation. Their inclusion must be stated before inter-
preting the residual.
Stored-component reductions reproduce the declared residuals. They test the arithmetic iden-
tity from saved components, not an independent derivation of those components from particle states
or reaction histories. A ledger can be arithmetically consistent while RE remains large. The base-
line TAF1 recorded boundary work is 1677.72 J m\ensuremath{-}1 after division by the numerical extrusion. Its
boundary projection term Db/bn is 674.26 J m\ensuremath{-}1, and its plane-strain term Dp/bn is 830.75 J m\ensuremath{-}1.
These terms are substantial on the same scale as the imposed work and cannot be omitted or
relabelled as physical fracture dissipation. They are per-unit-thickness numerical quantities; the
experimental thickness remains unspecified here.
Table 2 shows a monotone reduction of the declared final energy defect across all five saved
levels. The baseline TAF2 final value is 87.6280\%. The first three TAF1 values reproduce the
originally summarized sequence; the two finer levels extend the descriptive trend. Independent
reductions of the stored component ledgers support the displayed values, including the two finer
levels. This is evidence that the retained temporal schedules affect the diagnostic strongly, while
the finest retained result still has a 6.9786\% defect.
There is a source-version qualification: the baseline and refined mechanics cores differ. Their
inspected text difference is restricted to exposure and validation of the CFL argument, scheduling
substitution and returned metadata. This narrows the change but is not a formal demonstration
of compiled or runtime equivalence. Consequently we do not fit an asymptotic order, extrapolate a
4

Table 2: TAF1 retained temporal ladder at fixed h = 1 mm , \ensuremath{\ell} = 2 mm and common final time
29.976 \ensuremath{\mu}s. The event is the first saved particle maximum d > 10\ensuremath{-}3. The first row uses the baseline
source; the remaining rows use the source exposing an explicit CFL parameter.
Mechanics CFL Steps Final 100\ensuremath{\eta}E Threshold time ( \ensuremath{\mu}s)
0.250000 324 88.3429 1.20274
0.125000 642 44.0329 1.30736
0.062500 1284 21.7618 1.47079
0.031250 2568 11.4524 1.49413
0.015625 5136 6.9786 1.43577
zero-defect limit or attribute the full sequence to timestep alone under a proven fixed executable.
Temporal reduction also leaves spatial, particle-density and grid-phase effects unresolved.

\section*{4 Field-reference and initial-state results}

Table 3: Global discrepancies on the seven common times. Percentages use Eq. ( 4); N is the number
of nodes in the specified mask. Each displacement reference has its own target denominator.
Specimen Mask N 100\ensuremath{\varepsilon}u,abs 100\ensuremath{\varepsilon}u,inc 100\ensuremath{\varepsilon}v
TAF1 All supported 3429 41.9292 42.0446 36.4432
TAF1 Interior 3360 43.7037 43.9188 37.6298
TAF2 All supported 3465 64.1526 54.0514 52.6399
TAF2 Interior 3398 65.9291 55.8402 54.0600
The original approximately 64\% displacement result for TAF2 corresponds to the absolute
target, while approximately 54\% corresponds to the incremental target (Table 3). The predic-
tion is unchanged. Both the error numerator and the target denominator change under reference
subtraction. Selecting the smaller percentage therefore cannot establish a numerical or physical
improvement. TAF1 provides a useful contrast: its relative displacement discrepancy changes only
slightly and increases under subtraction.
Dimensional discrepancies retain the same distinction. The all-node displacement vector RMSE
is 60.84 \ensuremath{\mu}m for the TAF1 absolute target and 58.41 \ensuremath{\mu}m for its incremental target; TAF2 gives
96.16 and 74.70 \ensuremath{\mu}m. TAF1’s smaller incremental RMSE alongside a slightly larger relative dis-
crepancy illustrates why the target denominator matters. Velocity RMSE is unchanged at 3.353
and 4.693 m s\ensuremath{-}1 for TAF1 and TAF2. Excluding the impacted-edge band increases each relative
discrepancy moderately, so the all-node aggregates are not concealing a close interior match.
The recovered initializer arrays have exactly zero initial velocity for both runs. Under the
retained private analysis mapping and its unit labels, the first measured velocity vector RMS is
3.180814 and 5.277886 m s\ensuremath{-}1. Those are nodewise vector RMS values, not projectile velocities. The
initial simulated velocity is zero, hence the first-frame velocity relative discrepancy is one. The
initial absolute measured displacement is also nonzero, whereas the simulated displacement is zero.
Reference subtraction removes this displacement offset but does not change the initial velocity
mismatch.
The private extractor separates the synthetic reference output from the nonzero-numbered
5

data outputs, and maps the first actual data output to the first time-vector entry, labelled zero.
The publication-linked exporter can prepend an artificial zero field, so a zero reference output
is not evidence of physical rest. This resolves the bookkeeping convention used for the retained
calculation, while leaving the physical initial-state assumption open. A complete camera-to-solver
clock join, upstream unit conversion, filtering and derivative provenance, and consistency of initial
displacement, velocity and stress remain necessary for a physical interpretation. Nonzero first-
retained velocity by itself cannot isolate a missing constitutive mechanism.

\section*{5 Damage threshold versus crack initiation}

The numerical event in Table 2 is the first saved time when any particle damage exceeds 10\ensuremath{-}3. TAF1
crosses between 1.20274 and 1.49413 \ensuremath{\mu}s over the retained ladder, without a monotone threshold
time. Baseline TAF2 crosses at 0.277556 \ensuremath{\mu}s; the TAF1 range must not be assigned to both speci-
mens. These values depend on a declared damage threshold and sampled solver clock.
Equation ( 3) explains why this is a weak crack indicator. For spatially homogeneous H, van-
ishing damage gradient and a zero lower obstacle, stationarity gives
dhom = 2(1 \ensuremath{-} k)H
Gc/\ensuremath{\ell} + 2(1 \ensuremath{-} k)H . (9)
Thus positive tensile history produces small diffuse damage even without a resolved crack surface
or opening. Equation ( 9) is a homogeneous subproblem relation; it does not reconstruct the hetero-
geneous specimen or predict the location of a crack. It does show why identifying a small threshold
crossing with crack initiation is unwarranted.
The experimental event has a different definition and clock. Eid and colleagues identify initiation
and branching from crack-lip opening at twin nodes along the observed path [ 1]. Their TAF1
account places an earlier crack to the left of the hole near 40 \ensuremath{\mu}s and the right-notch crack near
60 \ensuremath{\mu}s after camera trigger. The latter is not the first crack anywhere in the specimen. First
retained DIC-derived elastic-parameter samples near 59.95 \ensuremath{\mu}s for TAF1 and 62.89 \ensuremath{\mu}s for TAF2 do
not, by file position alone, establish identical event criteria for both cases.
The camera trigger includes an experimental timing convention associated with wave travel
through the waveguide. Its correspondence to the solver’s initial state has not been fully established.
Moreover, the paired field comparison ends at 29.976 \ensuremath{\mu}s, before those reported TAF1 opening
events. No unqualified numerical-to-experimental initiation-time ratio follows from these records.
A defensible timing comparison would first define a simulated opening or path event comparable
to the observed one, assign its spatial location, and join the two clocks.

\section*{6 Interpretation and limits}

Three results constrain the interpretation of M0. First, the energy diagnostic is strongly schedule-
sensitive and remains finite at the finest retained level. Its projection terms are explicitly numerical
bookkeeping. Second, a substantial field mismatch persists under both displacement conventions
and in the interior mask. The material response is not the only uncontrolled cause because the
imposed boundary, reference subtraction and initial state enter the same comparison. Third, a
small AT2 crossing is not equivalent to an experimentally resolved crack-opening event. Treating
these three diagnostics as interchangeable evidence of a failed fracture law would exceed what they
measure.
6

The retained evidence therefore does not support using this setup as a validated high-rate
predictor. It also does not establish that all rate-independent AT2 models fail for the experiment,
or that adding rate dependence would resolve the discrepancies. Such attribution would require a
controlled comparison after numerical and input consistency have been addressed. The exploratory
Gc and \ensuremath{\ell}, inherited effective elasticity, at-rest initialization and constraint treatment remain coupled
choices in M0.
Several limits remain material to the result. The analysis reproduces saved-array differences
and stored-ledger arithmetic, rather than independent particle-energy, reaction-force or projection
replay. There is no established spatial, particles-per-cell, grid-phase or length-scale convergence
study for this comparison, no demonstrated zero-defect temporal limit, and no complete executed-
runtime equivalence proof. The experimental processing record does not yet supply a complete
velocity uncertainty model or a full initial-state and clock reconciliation. Repeat-run evidence is
not established by these two specimen records. Calibration dependence precludes independent
physical validation, and unknown experimental thickness prevents absolute experimental three-
dimensional energy claims. A separate unresolved-evidence report specifies the missing checks;
these limitations are part of the scientific scope, not additional measured failures.

\section*{7 Conclusions}

The PMMA case supports a bounded diagnosis of a coupled APIC/MPM–AT2 setup. Its five
retained TAF1 schedules reduce the declared final energy residual from 88.3429\% to 6.9786\%,
with source-version and convergence qualifications. Its displacement comparison depends materi-
ally on the reference convention for TAF2, while large velocity and interior discrepancies remain.
Its first d > 10\ensuremath{-}3 event is a damage-threshold crossing and cannot be substituted for observed
crack initiation. These results preserve the negative evidence while narrowing its inference: nu-
merical balance, model inputs and observable definitions must be resolved before a fracture-law or
predictive-accuracy claim can be assessed.
Evidence availability . The retained computational evidence is a private source bundle, desig-
nated Supplement S1, containing the source bindings, saved-result definitions and reduction records
underlying this audit. It is not a published dataset or an independently replayed solver result. The
experimental article and accompanying public data are cited below. No new experimental data are
claimed.

\section*{References}

[1] E. Eid, R. Seghir and J. Réthoré, “Crack branching at low tip speeds: spilling the T,” Journal
of Theoretical, Computational and Applied Mechanics (2023). doi:10.46298/jtcam.10172.
[2] E. Eid, R. Seghir and J. Réthoré, accompanying experimental data for “Crack branching at low
tip speeds: spilling the T,” Zenodo (2023). doi:10.5281/zenodo.8256346.
7

\end{document}
